\section{Rate-Aware Protection Assignment}\label{sec:method}
The cost $\hat J_p$ of Eq.~\eqref{eq:jhat} depends on $p$ and on the residual curves of the available mechanisms, so minimizing it gives an assignment that changes with the fault rate by construction. We minimize it for a given $p$ and storage budget $S$.

\paragraph{Ladders and moves.} The mechanisms available to a bit form a ladder of levels with equal storage steps: none and ideal suppression (one step of one bit), or none, three copies and five copies (two steps of two bits). Every move changes one bit by one level, so swapping $m$ step-ups for $m$ step-downs keeps the storage exactly at $S$. Codes and checksum correction protect whole words or channels and are compared as fixed, blanket mechanisms.

\paragraph{Value of a move.} The first-order change of $\hat J_p$ when bit $b$ moves from level $s$ to $s'$, per storage bit, is $-v_b$ with
\begin{equation}\label{eq:value}
v_b=\frac{(\rho_s-\rho_{s'})\big[\hat g_{w(b)}\Delta_b+\tfrac12(1-\rho_s-\rho_{s'})F_{w(b)}\Delta_b^2\big]}{\mathrm{cost}_{s'}-\mathrm{cost}_s}.
\end{equation}
The gradient $\hat g$ is taken at the average faulty weights $\theta+\mu$ of the current assignment, not at $\theta$. At small $p$, $\mu\approx0$ and $v_b$ is the signed per-bit term of Eq.~\eqref{eq:p-expansion}. At larger $p$ the gradient at $\theta+\mu$ reflects the overshoot of the current remainder, so bits whose flips would push further in the same direction gain value and bits whose flips push back lose it. The switch of Section~\ref{sec:theory} happens inside the search, without a threshold.

\paragraph{Search.} Starting from the better of two assignments, the absolute-gradient ranking and the ranking by $v_b$ at $\mu=0$, each round computes $\mu$ and $\hat g$ with one forward and one backward pass over the calibration inputs, swaps the $m$ best step-ups for the $m$ cheapest step-downs, and keeps the swap if $\hat J_p$ falls; otherwise $m$ is halved. The search stops when $m<16$ or after 40 rounds. Every round records the cost and its parts, the storage used, the number of bits at each level and the distribution of $v_b$ over all moves and over the chosen ones, so the course of the optimization can be inspected (Appendix~\ref{app:runr}).

\paragraph{Calibration noise.} Each mean gradient is shrunk by its standard error over the calibration inputs, $\tilde g_w=\bar g_w\max(0,1-\mathrm{se}_w^2/\bar g_w^2)$, so the signed term vanishes where its sign is not resolved.

\paragraph{Monte Carlo correction.} The part of $\hat J_p$ that is approximated is the spread around the average weights. With $K$ fault fields drawn once and shared by every candidate assignment, the correction $\hat C=\frac1K\sum_k[L(\theta+X_k)-L(\theta+\mu)]-\tfrac12\sum_wF_wV_w$ replaces that approximation by a sample, and the mean gradient over the $K$ faulted models replaces $\hat g$. The correction samples only what the expression misses, and the shared fields make differences between candidates far less noisy than the objective itself. \TBD{$K$ and the refresh interval, from the Stage 0 rule.}

\paragraph{Annealing.} As an optional variant, swaps are sampled with probability proportional to $\exp(v_b/T)$ and a worse assignment is accepted with probability $\exp(-\Delta J/T)$, with $T$ set by the first rejected step and lowered by 10\% per round. It is kept only if it improves on plain descent (rule R2c).

\paragraph{Outputs.} For a rate or a range of rates and a set of available mechanisms, the method returns the assignment, the predicted loss of every candidate against $p$, the predicted switch rates between the simple rankings, the loss of each assignment when the rate is misjudged, and a map of how many bits of each layer and bit position receive each mechanism. A designer who knows the operating regime, for example from the scrubbing interval, can read off whether a simple ranking is enough there.

\paragraph{Cost.} On ResNet-8 one analytic assignment takes \TBD{seconds, from Stage 0} on two CPU threads, and the corrected one \TBD{seconds}. The per-input gradients are those already needed by the absolute-gradient ranking.
